% TOP_PROBLEM: {"id": 3237, "title": "Bounding the chromatic number of triangle-free graphs with fixed maximum degree", "category_id": 3, "category": {"id": 3, "name": "graph_theory", "display_name": "Graph Theory", "description": "Problems involving graphs, networks, and their properties.", "slug": "graph-theory", "order_index": 3, "created_at": "2026-07-31T15:26:25.671Z"}, "status": "open", "problem_number": "OPG-36907", "set_id": 12, "statusQualification": "The source ceiling formula is preserved; the stronger floor version mentioned in its discussion is a different target."}
% FIRST_POSTED: 2026-10-02
% ATTEMPT_STATUS: unresolved
% UnsolvedMath ID 3237; source catalogue code OPG-36907
% !TeX program = lualatex
% Research attempt by GPT-6 Astra Ultra; no claim of a new complete solution.
\documentclass[11pt,a4paper]{article}
\usepackage[margin=25mm]{geometry}
\usepackage{fontspec}
\setmainfont{lmroman10-regular.otf}[BoldFont=lmroman10-bold.otf,
 ItalicFont=lmroman10-italic.otf,BoldItalicFont=lmroman10-bolditalic.otf]
\usepackage{amsmath,amssymb,mathtools}
\usepackage{unicode-math}
\setmathfont{latinmodern-math.otf}
\AtBeginDocument{\let\setminus\smallsetminus}
\usepackage{xurl,hyperref}
\hypersetup{colorlinks=true,urlcolor=blue}
\setcounter{secnumdepth}{0}
\setlength{\parindent}{0pt}
\setlength{\parskip}{5pt}
\setlength{\emergencystretch}{3em}
\begin{document}
\section*{Bounding the chromatic number of triangle-free graphs with fixed maximum degree}\label{3237}
{\small UnsolvedMath ID 3237 · OPG-36907 · Graph Theory · OpenGarden}
\subsection{Definitions and mathematical statement}
A finite simple undirected graph $G=(V,E)$ is triangle-free if no three
vertices are pairwise adjacent. Its maximum degree is
$\Delta(G)=\max_{v\in V}|\{u:uv\in E\}|$. A proper $q$-colouring is a map
$c:V\to\{1,\ldots,q\}$ with different values on adjacent vertices, and
$\chi(G)$ is the least possible $q$. For the empty graph set
$\chi(G)=\Delta(G)=0$.

The precise inequality displayed in this catalogue record is the following:
for every finite triangle-free simple graph,
\[
                \chi(G)\le\left\lceil\frac{\Delta(G)}2\right\rceil+2.
\]
The ceiling is the least integer greater than or equal to its argument.
Equivalently, for every integer $D\ge0$, every triangle-free graph of
maximum degree at most $D$ should admit a proper colouring with
$\lceil D/2\rceil+2$ colours. The colour bound is independent of the number
of vertices and does not assume regularity or connectedness.

There is a formulation issue worth making explicit. The source discussion
relates this to Reed's bound involving clique number, whose triangle-free
specialization gives the stronger $\lfloor\Delta/2\rfloor+2$ bound. The
present definition preserves the displayed ceiling version; for odd
$\Delta$ the two targets differ by one colour. For example, at $\Delta=6$
the displayed question asks for five colours, while at $\Delta=7$ it asks
for six. Triangle-free excludes only cycles of length three, so cycles of
length four and longer remain permitted.
\subsection{Short English statement}
Can every triangle-free graph be properly coloured with at most half its maximum degree, rounded up, plus two colours?
\subsection{Research attempt}
\paragraph{Scope and process.}
This is a GPT-6 Astra Ultra research attempt dated 2 October 2026, using
primary-source browsing and elementary mathematical checks. The canonical
definition above is retained. Results below are restricted results or
reductions, not a claim of a new solution to the entire catalogue question.
No formal proof assistant or external mathematical referee checked this text.


\paragraph{Literature and the exact rounding convention.}
Kang and Rosenfeld [R1] prove the stronger bound
$\chi(G)\le\lceil(\Delta+1)/2\rceil+1$ for triangle-free graphs
with $\Delta\ge524$. For even $\Delta$ this equals the catalogue
bound; for odd $\Delta$ it saves one additional colour. Thus the
large-degree range is already covered by that external theorem. We do
not reprove it. The independent work here gives a critical-graph
restriction and verifies an infinite family with unbounded chromatic
number, making explicit why triangle-freeness alone cannot reduce the
question to two-colouring. All graphs and colourings below are ordinary,
not list-coloured.

\paragraph{An elementary counterexample reduction.}
The greedy bound $\chi\le\Delta+1$ settles $0\le\Delta\le3$,
including the empty graph separately. If the displayed target fails,
choose a counterexample with the fewest vertices and put
$q=\lceil\Delta(G)/2\rceil+2$. Every graph obtained by deleting one
vertex satisfies the proposed inequality and hence is $q$-colourable,
since its maximum degree cannot increase. If a vertex $v$ had fewer
than $q$ neighbours, a $q$-colouring of $G-v$ could be extended to
$v$ by an unused colour. Therefore every such minimal counterexample
satisfies
\[
 \delta(G)\ge\lceil\Delta(G)/2\rceil+2.
\]
It is also connected: if every component satisfied its own bound,
reusing one common palette across components would colour $G$ within
its bound. Combined with [R1], this confines the missing degree range
to $4\le\Delta\le523$. This is a finite range of degrees, not a
finite collection of graphs: the order remains unbounded.

\paragraph{A fully analysed family with arbitrarily many colours.}
For a graph $H$ with vertices $v_1,\ldots,v_n$, construct $M(H)$ by
retaining $H$, adding independent shadow vertices $u_1,\ldots,u_n$,
joining $u_i$ to every original neighbour of $v_i$, and adding a
vertex $w$ adjacent to every shadow. If $H$ is triangle-free, so
is $M(H)$. A triangle through $w$ would require an edge between
shadows. A triangle through exactly one shadow $u_i$ and two original
vertices would give a triangle through $v_i$ in $H$. There is no
triangle through two shadows, and original triangles were excluded.

If $\chi(H)=k$, colour $v_i$ and $u_i$ alike in a $k$-colouring
of $H$, then give $w$ a new colour. Thus $\chi(M(H))\le k+1$.
For the reverse inequality, suppose $M(H)$ has a $k$-colouring and
call the colour of $w$ colour $k$. No shadow has that colour.
Recolour every original vertex of colour $k$ with the colour of
its shadow. Two adjacent original vertices were not both colour $k$,
and a recoloured $v_i$ cannot conflict with an unchanged neighbour
$v_j$, because $u_i$ is adjacent to $v_j$. This gives a proper
$(k-1)$-colouring of $H$, a contradiction. Hence
$\chi(M(H))=\chi(H)+1$.

Start with $H_0=C_5$ and set $H_{t+1}=M(H_t)$. Direct degree
counting gives
\[
 |V(H_t)|=6\cdot2^t-1,\qquad \chi(H_t)=t+3,
 \qquad \Delta(H_t)=3\cdot2^t-1\quad(t\ge1).
\]
Indeed original degrees double, shadow degrees become $d_H(v)+1$,
and the apex has degree $|V(H)|$; these give the asserted recurrence
and its base case. For $t\ge1$, the catalogue bound becomes
$3\cdot2^{t-1}+2$, which is at least $t+3$ by elementary induction.
Thus every member satisfies the target, despite unbounded chromatic
number. The construction tests the interaction of degree and chromatic
number; it is not an extremal family for the proposed linear bound.

\paragraph{Verification and remaining obstruction.}
The companion script
\texttt{scripts/check\_attempt\_batch03\_colouring\_2026\_10\_02.py}
constructs the first four members, checks their orders, degrees and
absence of triangles, and checks the inductively supplied colourings.
The chromatic lower bound is the proof above, not a numerical search.
A minimal-counterexample reduction controls local degrees but supplies
no recolouring move when every vertex sees all $q$ colours. The
Mycielski argument also shows that avoiding triangles does not prevent
global colouring obstructions.

\paragraph{Final status.}
\textbf{UNRESOLVED by this attempt.} The cited theorem settles large
degrees, and the independent restrictions and infinite family are
verified. A uniform argument for the remaining finite degree range,
with arbitrarily large order and the exact ceiling convention, is missing.

\subsection{Sources}
\begin{itemize}
\item[S1] ULAM.AI and the UnsolvedMath Contributors, supplied problems.json, record 3237 (OPG-36907), OpenGarden collection. Catalogue identity and source transcription; CC BY 4.0.
\url{https://huggingface.co/datasets/ulamai/UnsolvedMath}.
\item[S2] Open Problem Garden, Bounding the chromatic number of triangle-free graphs with fixed maximum degree. Original problem and discussion; the mathematical formulation below is expanded from the supplied source record.
\url{http://www.openproblemgarden.org/op/bounding_the_chromatic_number_of_triangle_free_graphs_with_fixed_maximum_degree}.

\item[R1] Ross J. Kang and Matthieu Rosenfeld, \emph{On Vizing's
problem for triangle-free graphs}, Electronic Journal of Combinatorics
32(4) (2025), P4.3. Primary theorem and rounding checked 2 October 2026.
\url{https://www.combinatorics.org/ojs/index.php/eljc/article/download/v32i4p3/pdf/}.

\end{itemize}
\end{document}
